\chapter{Literature Survey}\label{chap:literature}

\littable{Patrick Lewis et al. \cite{lewis2020rag}}
  {Retrieval-Augmented Generation for Knowledge-Intensive NLP Tasks}{2020}
  {Pretrained BART generator combined with a dense vector index of Wikipedia (DPR retriever,
   FAISS)\newline RAG-Sequence and RAG-Token variants\newline End-to-end training without labels
   for the correct passage}
  {Natural Questions, TriviaQA, WebQuestions, CuratedTrec\newline MS MARCO\newline Jeopardy
   question generation\newline FEVER}

\littable{Nelson F. Liu et al. \cite{liu2024lost}}
  {Lost in the Middle: How Language Models Use Long Contexts}{2024}
  {Controlled experiments that move the relevant information to different positions in the
   input\newline Multi-document QA with 10, 20 or 30 passages\newline Synthetic key-value
   retrieval}
  {NaturalQuestions-Open\newline Synthetic JSON key-value pairs}

\littable{Shahul Es; Jithin James; Luis Espinosa Anke; Steven Schockaert \cite{es2024ragas}}
  {RAGAs: Automated Evaluation of Retrieval Augmented Generation}{2024}
  {Reference-free metrics computed by prompting an LLM: faithfulness, answer relevance and context
   relevance}
  {WikiEval: built by the authors from 50 Wikipedia pages, with human judgements}

\littable{Akari Asai et al. \cite{asai2024selfrag}}
  {Self-RAG: Learning to Retrieve, Generate, and Critique through Self-Reflection}{2024}
  {Reflection tokens (Retrieve, IsRel, IsSup, IsUse)\newline Critic trained on GPT-4
   labels\newline Fine-tuned Llama 2 generator}
  {PopQA, TriviaQA\newline PubHealth\newline ARC-Challenge\newline Biographies, ASQA}

\littable{Shi-Qi Yan et al. \cite{yan2024crag}}
  {Corrective Retrieval Augmented Generation}{2024}
  {Fine-tuned T5 retrieval evaluator with three confidence levels\newline Refinement of the
   retrieved documents\newline Web search as fallback}
  {PopQA\newline Biography\newline PubHealth\newline ARC-Challenge}

\littable{Cheng Niu et al. \cite{niu2024ragtruth}}
  {RAGTruth: A Hallucination Corpus for Developing Trustworthy Retrieval-Augmented Language
   Models}{2024}
  {About 18,000 responses of six LLMs with word-level human labels\newline Four hallucination
   types (evident/subtle conflict, evident/subtle baseless information)\newline Comparison of
   detection methods}
  {MS MARCO (QA)\newline Yelp business data (data-to-text)\newline News summarisation}

\section{Related Work in Detail}

\subsection{Retrieval-Augmented Generation}
Lewis et al.\ \cite{lewis2020rag} combine a pretrained generator with a dense index of Wikipedia
and report results above much larger closed-book models on open-domain QA benchmarks. In their
human evaluation, generations with retrieval were judged more factual and more specific than
those without. Replacing the index changed the answers about world leaders, so knowledge can be
updated without retraining. The paper does not directly measure whether answers stay faithful to
the retrieved passages. That question is the starting point of this thesis.

Liu et al.\ \cite{liu2024lost} show that models do not use all of their input equally well.
Accuracy follows a U-shaped curve: it is highest when the relevant passage is at the start or end
of the input and drops clearly when it sits in the middle. Beyond a certain point, more passages
barely improve accuracy. For this thesis, these findings motivate a small fixed number of
passages, and they support passing passages in the order given by the retriever.

\subsection{Measuring Faithfulness and Hallucination}
Es et al.\ \cite{es2024ragas} propose metrics for evaluating RAG without reference answers. Their
faithfulness metric asks an LLM to split an answer into statements and to check each statement
against the context. Its scores agreed with human judgements better than simple baselines, but the
judge is itself an LLM and can make mistakes. The judge in this thesis is adapted from this method
(Section~\ref{sec:judge}), and it is checked against manual labels for this reason.

Niu et al.\ \cite{niu2024ragtruth} built a large corpus of RAG outputs with human hallucination
labels. Every model hallucinated, with rates that differed between models and tasks, and prompting
a strong LLM to detect hallucinations worked only moderately well. Their distinction between
conflicts with the context and baseless information is the basis of the labels used here.

\subsection{Reducing Hallucination in RAG}
Self-RAG \cite{asai2024selfrag} trains a model to generate reflection tokens that decide when to
retrieve and whether an output is supported. It reports clear gains in factuality and citation
accuracy, but it requires fine-tuning, so it cannot be applied to models that are only available
through an API. CRAG \cite{yan2024crag} adds a small fine-tuned evaluator that scores the retrieved
documents, and it refines or replaces them with web search depending on the score. Its results show
that checking retrieval quality before generation makes RAG more robust, especially when retrieval
is poor. The evaluator must still be trained, and web search adds latency and cost.

\subsection{Research Gap}
Self-RAG and CRAG show that deciding when to trust retrieval can reduce hallucination, but both
depend on a fine-tuned component. This thesis studies a simpler variant: a retrieval score, or the
share of the answer's words found in the passages, serves as the confidence signal, and the system
withholds the answer instead of correcting it. The rule needs no training or change to the generator,
so it works with any model behind an API. The thesis also measures what the rule costs in correct
answers, on answers not used to choose its threshold.

Recent work is close to this in aim. HALT-RAG \cite{goswami2025halt} combines NLI models and lexical
signals in a trained, calibrated meta-classifier that supports abstention. Bouke \cite{bouke2026context}
proposes a training-free detector that compares a sentence's likelihood with and without each source
chunk, and finds it weakest on short-answer QA. Sinha \cite{sinha2025semantic} reports that
embedding-based detectors fail on real hallucinations from aligned models. These papers come from a
short web search on 6 November 2026, not a systematic review, so this thesis does not claim that such
signals are new. Its aim is narrower: to compare very cheap signals that any pipeline already has, and
to relate them to each generator's errors on multi-hop questions.
